\section{Solution method and verification}
\label{sec:numerics}

\subsection{Thickness discretisation}
\label{sec:numerics_basis}

The thickness dependence is represented by a spectral Galerkin expansion.
Because the enthalpy contains the strain gradient, the weak form contains second
derivatives through the thickness and the trial space must be at least $C^1$ in
$z$.  Legendre polynomials are used for all fields, in the scaled variable
$\zeta=2z/d\in[-1,1]$,
\begin{equation}
\label{eq:basis}
U_i(x,\zeta)=\sum_{j=0}^{n-1}U_{ij}(x)\,P_j(\zeta),\qquad
N(x,\zeta)=\sum_{j=0}^{n-1}N_{j}(x)\,P_j(\zeta),
\end{equation}
with $P_j$ the Legendre polynomial of degree $j$ and $n$ the number of
thickness modes.  The two electrical cases use different potential spaces,
which is how the essential and natural conditions of
\cref{sec:governing_bc} are realised:
\begin{itemize}
\item grounded, \eqref{eq:bc_grounded}: the potential basis is formed from the
combination $P_j-\sqrt{(2j+1)/(2j+5)}\,P_{j+2}$, which vanishes at
$\zeta=\pm1$ for every $j$, so $\varphi=0$ holds on both faces exactly and the
space has $n-2$ members;
\item open circuit, \eqref{eq:bc_open}: the constant mode is removed and the
remaining $n-1$ modes are unconstrained, so $D_z=0$ is imposed weakly through
the surface term of the weak form.
\end{itemize}
The mechanical conditions \eqref{eq:bc_faces_mech} and the blocking condition
\eqref{eq:bc_faces_carrier} are natural and appear as the corresponding boundary
rows of the discretised operator.  The basis is integrated with
$\max(64,2n+10)$ Gauss--Legendre points, which resolves the mass and stiffness
matrices to machine precision for every degree used here.

The semi-discrete system is assembled as a three-term polynomial in the
propagation wavenumber $k_x$,
\begin{equation}
\label{eq:polynomial}
\Bigl[\sum_{r=0}^{2}\sum_{s=0}^{2}(-\ii k_x)^{r}(\ii k_x)^{s}\,
\mathsf{A}_{rs}(\omega,k_y)\Bigr]\vect{q}=0 ,
\end{equation}
where $\vect{q}$ collects all thickness degrees of freedom and the matrices
$\mathsf{A}_{rs}$ contain the six blocks of \cref{tab:conjugates}.  The degree
in $k_x$ is four, which is the order of the hyperelastic operator; the classical
limit $\ell=0$ returns the familiar quadratic problem.

\subsection{Temporal and spatial routes}
\label{sec:numerics_routes}

Two independent solution routes are implemented.

The \emph{temporal} route fixes a real $k_x$ and solves the generalised
eigenvalue problem in $\omega^2$ obtained from \eqref{eq:polynomial} with the
potential and the electrochemical potential eliminated.  It produces the
dispersion branches $\omega(k_x)$, the group velocity
$v_g=c_s\,\partial\omega/\partial k_x$ by central differences along a
branch-tracked sequence, and the temporal attenuation
$\gamma=-2\Imag\omega$.  Branch tracking is done by a modal-assurance
criterion on the mass-orthogonal displacements, so a curve can be followed
through a crossing without jumping to a neighbouring branch.

The \emph{spatial} route fixes a real frequency and solves the quartic problem
\eqref{eq:polynomial} directly for $k_x$.  It is the route that is physically
relevant at an edge, because reflection is a process at fixed frequency, and it
is the only route that gives the spatial attenuation
$\alpha=2|\Imag k_x|$ used in \cref{fig:fig01,fig:fig06}.  Reporting the
imaginary part of $\sqrt{K^2+k_y^2}$ for a real $K$ as an attenuation would
return zero identically and is not done anywhere in this work.

Both routes are cross-checked against each other (\cref{fig:fig09}d) and against
strong-form integrations of the reduced system
(\cref{sec:numerics_verification}).  They are two implementations of the same
dispersion relation; agreement between them tests the implementation, not the
physics, and is described that way throughout.

\subsection{Edge matching}
\label{sec:numerics_edge}

At fixed $\omega$ the spatial problem returns $2M$ eigenvalues $k_x^{(j)}$ and
their states $\vect{y}^{(j)}$, with $M$ the number of unknowns in
\eqref{eq:polynomial}.  Each is classified as outgoing or incoming by

\begin{equation}
\label{eq:classification}
\text{outgoing}\iff
\Imag k_x^{(j)}<-\tau
\quad\text{or}\quad
\bigl|\Imag k_x^{(j)}\bigr|\le\tau\ \text{and}\ \mathcal{P}_x^{(j)}<0 ,
\end{equation}
where $\tau=2\times10^{-8}\max(1,|k_x|)$ and $\mathcal{P}_x$ is the
thickness-integrated power flux \eqref{eq:Px_int}.  A purely real root is
therefore outgoing when it carries energy away from the edge, and an
evanescent root is outgoing when it decays away from the edge.  The solver
verifies that the number of outgoing roots equals $M$ and aborts otherwise;
this is a structural check on the eigenvalue count, not a tolerance.

The solution is written as the incident mode plus a combination of the outgoing
modes,
\begin{equation}
\label{eq:edge_ansatz}
\vect{y}(0)=\vect{y}_{\text{inc}}+\sum_{j}a_j\,\vect{y}^{(j)},
\end{equation}
and the coefficients $a_j$ follow from the homogeneous edge conditions
\eqref{eq:bc_edge_mech}--\eqref{eq:bc_edge_carrier} written as
$\mathsf{B}\vect{a}=-\vect{b}_{\text{inc}}$.  The system is row- and
column-scaled and solved in the least-squares sense through a singular-value
decomposition.  The rank of $\mathsf{B}$ is required to equal its column count
with a cut-off of $10^{-12}$ of the largest singular value; a rank deficiency
raises an error instead of returning a regularised answer.  The reported
condition number at the operating point is $1.6\times10^{2}$, and the scaled
residual of \eqref{eq:edge_ansatz} is below $10^{-12}$ for every row used in
this paper.

The backward channel is the outgoing mode whose real wavenumber is closest to
the target $k_{Bx}$, and its share of the incident power is
\begin{equation}
\label{eq:RB}
R_B=\frac{-\,P_{B}\,e^{2\Imag k_{Bx}L}}{P_{\text{inc}}\,e^{2\Imag k_{ix}L}}
=\frac{-\,P_{B}}{P_{\text{inc}}}\,e^{2(\Imag k_{Bx}-\Imag k_{ix})L},
\end{equation}
evaluated at the reference plane $x=-L$, $L=10\,d=\SI{2}{\micro\metre}$.  The
two exponentials are retained explicitly so that the ratio is measured at the
same plane for both modes.  Equation \eqref{eq:RB} defines a \emph{modal share}
of the outgoing power.  It is not a reflectance, and the manuscript never
treats it as one: the quantity that would correspond to a reflectance is the
sum over all propagating outgoing modes, which is reported separately in
\cref{fig:fig05}.

\subsection{Admissibility of a reported value}
\label{sec:numerics_admissibility}

A solved row is accepted only if it satisfies three conditions, none of which
is a plotting threshold.

\emph{Conservation.}  The backward mode is one of the outgoing modes, so its
self-power cannot exceed the power available at the plane,
\begin{equation}
\label{eq:admissible}
0\le R_B\le1+\frac{|P_{\text{cross}}|}{P_{\text{inc}}} .
\end{equation}
Across the certified solves the interference term is between $10^{-6}$ and
$10^{-2}$ of the incident power, so \eqref{eq:admissible} is applied with a
bound of $1.01$.  Rows that violate it would imply that a single outgoing mode
carries more power than the incident wave plus the interference, which is
impossible for a solution of the boundary-value problem; they indicate that the
identified mode is not the physical continuation of the backward branch and are
discarded.  An earlier version of this work instead discarded rows whose energy
angle exceeded $80^{\circ}$.  That was a plotting convention, not a physical
criterion, and it is replaced here by \eqref{eq:admissible}.

\emph{Convergence.}  The reported quantity must be stable under refinement of
the Galerkin degree.  The refinement study in
\cref{sec:results_peak} shows that both the operating-point value and the peak
value change by less than $10^{-6}$ between degree 18 and degree 30.

\emph{Resolution of the measured projection.}  $R_B$ is a ratio of $x$ fluxes.
When the energy angle of a mode approaches $\pm90^{\circ}$ its $x$ flux
vanishes and the ratio loses meaning even though the mode itself is perfectly
well resolved.  Rather than exclude such rows silently, the solver stores the
flux fraction $|\cos\theta_E|=|P_x|/\sqrt{P_x^2+P_y^2}$ next to every row so
that the reader can see how well the projection is resolved.  The two values
quoted in this paper have $|\cos\theta_E|=0.88$ and $0.97$ respectively.

\subsection{Verification}
\label{sec:numerics_verification}

The implementation is verified in five independent ways.  All of the following
run as an automated suite (\texttt{paper\_a\_study/test\_model.py}, 18 tests,
all passing in \SI{8.7}{\second}).

\begin{enumerate}
\item \emph{Exact classical limit.}  With $\ell=f=H=0$ the plate supports
exact shear-horizontal modes $\omega=\sqrt{c_{66}K^2+(j\pi)^2/d^2}$; the spectral
solution reproduces the first three of them to $10^{-9}$.
\item \emph{Exact strain-gradient limit.}  With $\ell\neq0$ the fundamental
shear-horizontal branch has the closed form
$\omega=\sqrt{c_{66}K^2(1+\ell^2K^2)}$, reproduced to $10^{-10}$.
\item \emph{Independent strong-form propagator.}  An eight-state classical
sagittal system and a twelve-state gradient system are integrated directly as
ODEs, with no spectral matrices.  They reproduce the spectral frequencies to
$10^{-10}$ for both the grounded and the open-circuit problem, and with
nonzero $f$ and $H$.
\item \emph{Independent discretisation.}  A $C^1$ finite-element plate
formulation in the thickness direction converges monotonically to the spectral
edge result, with an error below $3\times10^{-5}$ on the finest mesh.  This is
the only check in the list that tests the edge solver rather than the
dispersion solver.
\item \emph{Residuals of the strong form.}  For every mode used, the momentum,
Gauss, carrier-continuity, traction, double-traction and blocking residuals are
below $10^{-8}$ in the normalised measure reported in
\cref{fig:fig09}(c), and the scaled edge residual is below $10^{-12}$.
\end{enumerate}

Two further checks are written as tests rather than as figures.  The
reciprocity of the operator at $+k_y$ and $-k_y$ holds to $10^{-13}$, which
verifies the symmetry of the discretised bilinear form including the $H$ term;
and a deliberately singular edge system is rejected with an explicit error
rather than being regularised into a plausible-looking number.

\subsection{Reproducibility}
\label{sec:numerics_repro}

Every plotted number is produced by two scripts and no value is transcribed by
hand.  \texttt{build\_figure\_data.py} solves the physics and writes the
CSV/JSON files; \texttt{generate\_final\_figures.py} reads those files and only
plots.  The refinement study of \cref{sec:results_peak} is produced by a third
script, \texttt{certify\_peak.py}.  The classical constants live in
\texttt{a0/materials.json}, which is included in the distributed package
because the model cannot start without it.  \Cref{app:repro} lists the exact
commands.
